\chapter{域扩张与 Galois 理论}\label{chap:galois-2}
现代``Galois 理论''研究域扩张及其自同构群. 它起源于十九世纪数学家提出的一个具体问题：是否所有``代数方程''都``可用根式求解''？换言之，每个 $\mathbb Q[X]$ 中的不可约多项式是否都有一个能用 $+,-,\times,\div$ 和 $\sqrt[n]{x}$ 表示的复数解？二次方程公式、Cardan 的三次方程公式和 Ferrari 的四次方程公式说明，次数不超过 $4$ 时可以做到；但 Galois（以及 Abel）给出了一个不能如此求解的五次多项式. 事实上，次数 $\ge5$ 时能如此求解的情形甚至很少. 为此，Galois 研究了多项式方程解的对称性集合（现在称为``Galois 群''），并发现一个多项式方程可用根式求解，当且仅当其对称群可解（这里使用当时尚不存在的现代群论意义）. 一个 $n$ 次方程的对称群嵌入对称群 $\mathfrak S_n$. 当 $n<5$ 时，$\mathfrak S_n$ 可解，这解释了经典公式的存在. 相反，$\mathfrak A_5$ 是单群，因而不可解；Galois 正是给出了一个对称群为 $\mathfrak A_5$ 的方程.
\section{域扩张的一般性质}\label{sec:galois-2-1}
\subsection{定义}\label{sec:galois-2-1-1}
\begin{definition*}
设 $k$ 是域. $k$ 的一个\newterm{galois-field-extension}{扩张}{field extension}是配备域同态 $k\longrightarrow K$ 的域 $K$.
\end{definition*}
注意，域同态就是环同态，因此 $k$ 的扩张无非是一个本身为域的 $k$-代数. ``扩张''这一名称的依据是域同态总是单射. 常常略去同态，简记为 $k\subset K$. 此时，$K$ 的\newterm{galois-subextension}{子扩张}{subextension}是 $K$ 中包含 $k$ 的子域 $K'$.
\paragraph{记号.} 设 $k\subset K$ 是域扩张，$\alpha\in K$. 记 $k[\alpha]$ 为 $\alpha$ 生成的 $K$ 的 $k$-子代数，$k(\alpha)$ 为 $\alpha$ 生成的 $K$ 的子扩张.
与往常一样，``生成''指``包含 $\alpha$ 的最小者''. 具体而言，$k[\alpha]$ 是将 $X$ 映到 $\alpha$ 的唯一 $k$-代数同态 $k[X]\longrightarrow K$ 的像，故作为 $k$-模由 $\alpha$ 的幂生成. 显然 $k[\alpha]\subset k(\alpha)$，所以 $k(\alpha)=\operatorname{Frac}(k[\alpha])$. 但 $k(\alpha)$ 并不总是某个同态 $k(X)\longrightarrow K$ 的像.
\begin{example*}
取 $k=\mathbb Q\subset K=\mathbb C$，$\alpha=i$. 此时 $\mathbb Q[i]=\mathbb Q\oplus\mathbb Qi$ 是域，因此 $\mathbb Q[i]=\mathbb Q(i)$ 在 $\mathbb Q$ 上的维数是 $2$，而 $\mathbb Q(X)$ 在 $\mathbb Q$ 上无限维. 域同态是单射，所以不存在域同态 $\mathbb Q(X)\longrightarrow\mathbb Q(i)$.
\end{example*}
\subsection{代数与超越的二分}\label{sec:galois-2-1-2}
\begin{proposition*}
设 $k\subset K$ 是域扩张，$\alpha\in K$. 记 $\varphi_\alpha:k[X]\longrightarrow K$ 为将 $X$ 映到 $\alpha$ 的 $k$-代数同态. 则有如下二分：
\begin{enumerate}[label=\roman*)]
\item $\varphi_\alpha$ 为单射，此时它诱导同构 $k[X]\xrightarrow{\sim}k[\alpha]$，并唯一延拓为同构 $k(X)\xrightarrow{\sim}k(\alpha)$. 特别地，$k[\alpha]$ 和 $k(\alpha)$ 在 $k$ 上都无限维.
\item $\varphi_\alpha$ 不是单射，此时有以下性质：
\begin{enumerate}[label=(\alph*)]
\item 它的核由唯一的首一不可约多项式 $f_\alpha\in k[X]$ 生成；
\item $\varphi_\alpha$ 诱导同构 $k[X]/(f_\alpha)\xrightarrow{\sim}k[\alpha]$；
\item $k[\alpha]$ 在 $k$ 上有限维，维数为 $\deg(f_\alpha)$；
\item $k(\alpha)=k[\alpha]$.
\end{enumerate}
\end{enumerate}
\end{proposition*}
\begin{proof}
在第一种情形中，只须证明 $\varphi_\alpha$ 唯一延拓为同构 $k(X)\xrightarrow{\sim}k(\alpha)$. 因为 $\varphi_\alpha$ 将每个非零 $f\in k[X]$ 映到 $K$ 中的可逆元，这由分式域的泛性质得到.
在第二种情形中，$\operatorname{Ker}(\varphi_\alpha)$ 由一个多项式生成，是因为 $k[X]$ 为主理想整环. 这个多项式在相差一个单位因子的意义下确定；乘以某个 $\lambda\in k^\times$ 可使其首一，由 $k[X]^\times=k^\times$，此时它唯一确定. 此外，由商的泛性质，$\varphi_\alpha$ 诱导同构 $k[X]/\operatorname{Ker}(\varphi_\alpha)\xrightarrow{\sim}k[\alpha]$；又因为 $k[\alpha]\subset K$ 是整环，$\operatorname{Ker}(\varphi_\alpha)$ 是素理想，故 $f_\alpha$ 不可约. 这证明了 (a)、(b). (c) 由第~\ref{sec:galois-1-4-3}~分目的第二个推论得到. 至于 (d)，只须证明 $k[\alpha]$ 是域. 有两种方法：回忆 $k[X]$ 的每个素理想都是极大理想，或使用下述本身也值得注意的引理.
\end{proof}
\begin{lemma*}
域 $k$ 上有限维的整代数 $A$ 是域.
\end{lemma*}
\begin{proof}
乘以 $x\in A\setminus\{0\}$ 是 $A$ 的单射 $k$-线性自同态，由秩定理，它是双射. 特别地，存在 $y$ 使 $xy=1$.
\end{proof}
\begin{definition*}
在命题的情形中，第一种情形称 $\alpha$ 在 $k$ 上\newterm{galois-transcendental-element}{超越}{transcendental}，第二种情形称它在 $k$ 上\newterm{galois-algebraic-element}{代数}{algebraic}. 后一种情形中的 $f_\alpha$ 称为 $\alpha$ 的\newterm{galois-minimal-polynomial-element}{极小多项式}{minimal polynomial}.
\end{definition*}
\begin{example*}
考虑扩张 $\mathbb Q\subset\mathbb C$. 复数 $i,j,\sqrt2$ 都在 $\mathbb Q$ 上代数. 最早构造出的超越数是``Liouville 数''，它们可以被有理数很好地逼近. 随后 Lindemann 证明，代数数 $a$ 所给出的 $e^a$ 都是超越数. 由 $e^{i\pi}=1$，这蕴含 $\pi$ 超越，顺带说明``化圆为方''不可能. 另一方面，我们仍不知道 $e^\pi$ 或 $\pi+e$ 这样的数是否超越.\footnote{译注：本段按原文保留. 关于 $e^a$ 的断言须排除 $a=0$；原式 $e^{i\pi}=1$ 的右端应为 $-1$；原文关于 $e^\pi$ 的未知性陈述亦有误.}
\end{example*}
\begin{remark*}
事实上，超越复数比代数复数``多得多''. 更准确地说，所有代数数组成的子集 $\overline{\mathbb Q}\subset\mathbb C$ 是可数的. 每个代数数都是某个整系数多项式的零点. 这些多项式与最终为零的整数序列一一对应，而这些序列组成可数集（练习！）. 因此，超越数组成的集合 $\mathbb C\setminus\overline{\mathbb Q}$ 不可数.
\end{remark*}
\subsection{代数无关性}\label{sec:galois-2-1-3}
设 $k\subset K$ 是域扩张，$(\alpha_i)_{i\in I}$ 是由集合 $I$ 编号的 $K$ 中元素族. 与上一分目一样，记 $k[(\alpha_i)_{i\in I}]$ 为这些 $\alpha_i$ 生成的 $K$ 的 $k$-子代数，$k((\alpha_i)_{i\in I})$ 为它们生成的 $K$ 的子扩张.
\begin{definition*}
若将每个 $X_i$ 映到 $\alpha_i$ 的 $k$-代数同态 $k[(X_i)_{i\in I}]\longrightarrow K$ 为单射，就称族 $(\alpha_i)_{i\in I}$ 在 $k$ 上\newterm{galois-algebraically-independent}{代数无关}{algebraically independent}.
\end{definition*}
当 $\alpha_i$ 代数无关时，定义中的同态唯一延拓为同构
\begin{align*}
k((X_i)_{i\in I}):=\operatorname{Frac}(k[(X_i)_{i\in I}])\xrightarrow{\sim}k((\alpha_i)_{i\in I}).
\end{align*}
\begin{example*}
如果在 $\mathbb C$ 中``随机''选取 $n$ 个元素，它们极有可能代数无关. 但对预先给定的数，证明其代数无关性十分困难；例如，我们不知道 $e$ 与 $\pi$ 是否代数无关. 人们猜想 Riemann 的 $\zeta$ 函数在奇数点的值 $\zeta(3),\zeta(5),\ldots$ 在 $\mathbb Q$ 上代数无关. Apéry 定理``仅仅''证明了 $\zeta(3)$ 的无理性，却已享有盛名，由此可见这一猜想的分量.
\end{example*}
\begin{remark*}
若 $I=I_1\sqcup I_2$（不交并），则 $k((\alpha_i)_{i\in I})=k((\alpha_i)_{i\in I_1})((\alpha_i)_{i\in I_2})$. 此外，$(\alpha_i)_{i\in I}$ 在 $k$ 上代数无关，当且仅当 $(\alpha_i)_{i\in I_1}$ 在 $k$ 上代数无关，且 $(\alpha_i)_{i\in I_2}$ 在 $k((\alpha_i)_{i\in I_1})$ 上代数无关.
\end{remark*}
\subsection{各类扩张}\label{sec:galois-2-1-4}
\begin{definition*}
扩张 $k\subset K$ 称为：
\begin{itemize}
\item \newterm{galois-finite-extension}{有限扩张}{finite extension}，若 $K$ 作为 $k$-向量空间有限维. 此时记 $[K:k]:=\dim_k(K)$，称为 $K$ 在 $k$ 上的\newterm{galois-extension-degree}{次数}{degree of an extension}.
\item \newterm{galois-algebraic-extension}{代数扩张}{algebraic extension}，若每个 $\alpha\in K$ 都在 $k$ 上代数.
\item \newterm{galois-finitely-generated-extension}{有限生成扩张}{finitely generated extension}，若 $K$ 作为域扩张由有限个元素生成.
\item \newterm{galois-simple-extension}{单扩张}{simple extension}，若 $K$ 作为域扩张由一个元素生成.
\item \newterm{galois-purely-transcendental-extension}{纯超越扩张}{purely transcendental extension}，若 $K$ 由一个在 $k$ 上代数无关的族生成.
\end{itemize}
\end{definition*}
\begin{remark*}
由定义，有限扩张是代数扩张；一个扩张代数，当且仅当它是有限子扩张的并. 此外，代数扩张有限，当且仅当它有限生成. 最后，若 $\alpha$ 在 $k$ 上代数，则 $k(\alpha)$ 有限且 $[k(\alpha):k]=\deg(f_\alpha)$.
\end{remark*}
\begin{example*}
扩张 $k\subset k(X_1,\ldots,X_n)$ 有限生成且纯超越.
\end{example*}
\subsection{相对代数闭包}\label{sec:galois-2-1-5}
\begin{lemma*}
设 $k\subset k'\subset K$ 是域扩张塔.
\begin{enumerate}[label=\roman*)]
\item $K$ 在 $k$ 上有限，当且仅当 $K$ 在 $k'$ 上有限且 $k'$ 在 $k$ 上有限. 此时 $[K:k]=[K:k'][k':k]$.
\item $K$ 在 $k$ 上代数，当且仅当 $K$ 在 $k'$ 上代数且 $k'$ 在 $k$ 上代数.
\end{enumerate}
\end{lemma*}
\begin{proof}
第一条的等价性显然. 为证等式，令 $n=[K:k']$，$m=[k':k]$. 作为 $k'$-向量空间有 $K\simeq k'^n$，作为 $k$-向量空间有 $k'\simeq k^m$，因此作为 $k$-向量空间有 $K\simeq(k^m)^n=k^{mn}$. 具体地，若 $\alpha_1,\ldots,\alpha_n$ 是 $K$ 在 $k'$ 上的基，$\beta_1,\ldots,\beta_m$ 是 $k'$ 在 $k$ 上的基，则 $\{\alpha_i\beta_j:i=1,\ldots,n;\ j=1,\ldots,m\}$ 是 $K$ 在 $k$ 上的基.
第二条的正向蕴含显然. 反之，取 $\alpha\in K$，记它在 $k'$ 上的极小多项式为 $f_\alpha=X^n+a_1X^{n-1}+\cdots+a_n\in k'[X]$. 于是 $\alpha$ 在 $k(a_1,\ldots,a_n)$ 上代数. 各个 $a_i$ 在 $k$ 上代数，故用第一条归纳可知 $k(a_1,\ldots,a_n)$ 在 $k$ 上有限. 因此 $k(a_1,\ldots,a_n,\alpha)$ 在 $k$ 上有限，特别地，$\alpha$ 在 $k$ 上代数.
\end{proof}
下述命题表明，每个扩张都包含唯一的极大代数子扩张.
\begin{proposition*}
设 $k\subset K$ 是域扩张. $K$ 中所有在 $k$ 上代数的元素组成的集合 $K_{\mathrm{alg}}$ 是域，称为 $k$ 在 $K$ 中的\newterm{galois-relative-algebraic-closure}{相对代数闭包}{relative algebraic closure}.
\end{proposition*}
\begin{proof}
设 $\alpha,\beta\in K$ 在 $k$ 上代数. 则 $k(\alpha)$ 在 $k$ 上有限，而 $\beta$ 更在 $k(\alpha)$ 上代数，所以 $k(\alpha,\beta)$ 在 $k(\alpha)$ 上有限，从而在 $k$ 上有限. 特别地，$\alpha+\beta$ 和 $\alpha\beta$ 在 $k$ 上代数.
\end{proof}
\begin{example*}
前面引入的 $\overline{\mathbb Q}$ 是 $\mathbb Q$ 在 $\mathbb C$ 中的相对代数闭包.
\end{example*}
\subsection{超越基与超越次数}\label{sec:galois-2-1-6}
扩张 $k\subset K$ 总包含一个纯超越子扩张 $k\subset K'$，使 $K'\subset K$ 代数. 只须取一个按包含关系极大的代数无关族生成的扩张为 $K'$. 下例表明 $K'$ 远非唯一.
\begin{example*}
设 $C$ 是方程 $X^2=Y^3-1$ 定义的曲线. 它的多项式函数环是 $A=\mathbb C[X,Y]/(X^2-Y^3+1)$，有理函数域是 $K=\operatorname{Frac}(A)=\mathbb C(Y)[X]/(X^2-Y^3+1)$. 域 $K$ 并非纯超越，却是纯超越域 $\mathbb C(Y)$ 的二次扩张. 也可写成 $K=\mathbb C(X)[Y]/(Y^3-X^2-1)$，可见它是纯超越域 $\mathbb C(X)$ 的三次扩张.
\end{example*}
然而，这个例子中的纯超越子域都只有一个不定元. 这可推广如下.
\begin{theorem*}
设 $k\subset K$. $K$ 中在 $k$ 上代数无关的极大族都有相同基数. 它们称为 $K$ 在 $k$ 上的\newterm{galois-transcendence-basis}{超越基}{transcendence basis}，其基数称为扩张 $k\subset K$ 的\newterm{galois-transcendence-degree}{超越次数}{transcendence degree}，记作 $\deg.\operatorname{tr}.(K/k)$.
\end{theorem*}
\begin{proof}
这里只证明 $K$ 有有限极大代数无关族的情形. 设存在代数无关元素 $\alpha_1,\ldots,\alpha_n$，使 $K$ 在 $k(\alpha_1,\ldots,\alpha_n)$ 上代数. 设 $\beta_1,\ldots,\beta_m$ 是另一个代数无关族. 只须证明 $m\le n$. 我们将反复使用下述引理.
\end{proof}
\begin{lemma*}
设 $\alpha,\beta\in K$ 各自在子域 $k'$ 上超越，但在 $k'$ 上代数相关. 则 $\alpha$ 在 $k'(\beta)$ 上代数.
\end{lemma*}
\begin{proof}
存在不可约多项式 $f\in k'[X,Y]$ 使 $f(\alpha,\beta)=0$. 展开 $f=\sum_{r\in\mathbb N}g_r(Y)X^r$，其中 $g_r\in k'[Y]$. 因 $f\ne0$，$g_r$ 不全为零；因 $\beta$ 超越，$g_r(\beta)$ 也不全为零. 所以 $\alpha$ 是一个系数在 $k'(\beta)$ 中的非零多项式的根.
\end{proof}
\begin{proof}[定理证明续]
令 $I_0:=\{1,\ldots,n\}$. 因 $\beta_1$ 在 $k$ 上超越，集合
\begin{align*}
\{I\subset I_0:\ \beta_1\text{ 在 }k((\alpha_i)_{i\in I})\text{ 上超越}\}
\end{align*}
包含 $I=\varnothing$，因而非空. 在其中选取极大的 $I_1$. 因 $\beta_1$ 在 $k(\alpha_1,\ldots,\alpha_n)$ 上代数，有 $I_1\subsetneq I_0$. 对每个 $j\in I_0\setminus I_1$，$\beta_1$ 与 $\alpha_j$ 在 $k((\alpha_i)_{i\in I_1})$ 上代数相关，引理保证 $\alpha_j$ 在 $k(\beta_1)((\alpha_i)_{i\in I_1})$ 上代数. 所以 $K$ 在 $k(\beta_1)((\alpha_i)_{i\in I_1})$ 上代数.
特别地，$\beta_2$ 在 $k(\beta_1)((\alpha_i)_{i\in I_1})$ 上代数，但在 $k(\beta_1)$ 上超越. 因此存在极大的 $I_2\subsetneq I_1$，使 $\beta_2$ 在 $k(\beta_1)((\alpha_i)_{i\in I_2})$ 上超越；同上，$K$ 在 $k(\beta_1,\beta_2)((\alpha_i)_{i\in I_2})$ 上代数. 归纳得到 $I_0$ 的子集 $I_m$，使 $K$ 在 $k(\beta_1,\ldots,\beta_m)((\alpha_i)_{i\in I_m})$ 上代数. 因为 $I_0\supsetneq I_1\supsetneq\cdots\supsetneq I_m$ 严格递减，有 $0\le|I_m|\le n-m$，故 $m\le n$.
\end{proof}
\begin{example*}
每个 $\mathbb Q(X_1,\ldots,X_n)$ 都可数，故 $\mathbb C$ 在 $\mathbb Q$ 上的超越次数无限.
\end{example*}
\begin{remark*}
若代数子集 $V\subset\mathbb C^n$ 的 $\mathcal O(V)$ 为整环，则整数 $\deg.\operatorname{tr}.(\mathcal M(V)/\mathbb C)$ 起着维数的作用. 可以证明，它也等于每条极大素理想链 $\mathfrak p_0=\{0\}\subsetneq\mathfrak p_1\subsetneq\cdots\subsetneq\mathfrak p_n\subsetneq\mathcal O(V)$ 的长度.
\end{remark*}
\begin{lemma*}
设 $k\subset k'\subset K$ 是两个域扩张. $K$ 在 $k$ 上的超越次数有限，当且仅当 $k'$ 在 $k$ 上及 $K$ 在 $k'$ 上的超越次数都有限. 此时
\begin{align*}
\deg.\operatorname{tr}.(K/k)=\deg.\operatorname{tr}.(K/k')+\deg.\operatorname{tr}.(k'/k).
\end{align*}
\end{lemma*}
\begin{proof}显然.\end{proof}

